\section{דיון}

התוצאה המרכזית של העבודה אינה שה־HMM ``מנבא את השוק''. להפך: גם HMM, גם מודלי \tech{Supervised Learning} וגם כללי ייחוס פשוטים מתקשים להפיק אות עקבי לחיזוי היום הבא מתוך ארבעת המאפיינים שנבדקו. דווקא התוצאה הזו מחדדת את ההבדל בין שתי שאלות שונות: האם המודל מנבא היטב, והאם הוא מייצג היטב את מצב השוק.

\subsection{מדוע חיזוי יום קדימה נשאר קשה?}

חיזוי סימן התשואה של יום בודד הוא בעיה רועשת מאוד. התנודתיות של השוק, הטווח היומי והפעילות במסחר יכולים לתאר היטב את רמת הסיכון, אבל אינם בהכרח קובעים אם המחיר יעלה או ירד מחר. לכן ייתכן שמאפיין יהיה שימושי מאוד לזיהוי ``סביבה תנודתית'' בלי להיות שימושי לחיזוי הכיוון ביום הבא.

העובדה שגם \tech{Logistic Regression}, \tech{Random Forest} ו־\tech{HistGradientBoosting} אינם מציגים יתרון עקבי תומכת בפרשנות שהקושי אינו ייחודי ל־HMM. במילים אחרות, לא נראה שמודל סיווג אחר מצליח לחלץ מהתכונות הקיימות אות ברור שה־HMM מפספס.

בנוסף, DPA מושפע מאיזון המחלקות. אם בתקופה מסוימת יש יותר ימי עלייה, כלל פשוט שמנבא עלייה לעיתים קרובות יכול להשיג תוצאה מעל \pct{50}. לכן נעשה שימוש גם ב־\tech{Balanced Accuracy} וב־\tech{ROC-AUC}. ברוב הנכסים מדדים אלה נשארו קרובים ל־\pct{50}, ולכן אין בסיס לטעון שקיים אות חיזויי חזק.

\subsection{מה ה־HMM כן מוסיף?}

כאן נמצא הערך המרכזי של המודל. גם אם איננו יודעים לחזות בביטחון אם מחר יהיה יום חיובי או שלילי, עדיין שימושי לדעת אם אנחנו נמצאים כרגע בתקופה שקטה, בתקופה תנודתית או בסביבה דמוית לחץ. מידע כזה מתאר את \textbf{ההקשר} שבו התשואה הבאה תתרחש.

HMM עושה בדיוק זאת: הוא בונה ייצוג מפורש של מצב חבוי ושל דינמיקת המעבר בין מצבים. עבור \asset{SPY}, ארבעת המצבים נבדלים בו־זמנית בתשואה, תנודתיות, טווח יומי, \tech{drawdown} ומשך שהייה. לכן המודל אינו רק מחלק ימים לפי ``תנודתיות גבוהה או נמוכה'', אלא יוצר תיאור רב־ממדי של סביבת השוק.

נוסף לכך, וקטור ה־\tech{posterior} מספק מידע על מידת הביטחון של המודל. רוב הזמן ההסתברות מרוכזת במצב אחד, אך סביב גבולות בין מצבים היא מתפזרת יותר. כפי שהוסבר קודם, הקשר בין אנטרופיה לבין מעבר מצב הוא בחלקו מכני משום ששני המדדים נגזרים מאותו \tech{posterior}. לכן אין לראות בו אימות חיצוני. עם זאת, הוא כן מספק מדד פנימי שימושי לשאלה עד כמה ההחלטה על זהות המשטר חד־משמעית.

הראיות החיצוניות מגיעות משני מקורות אחרים. ראשית, \asset{VIX}, שלא שימש באימון, מקבל רמות שונות במצבים שונים. שנית, הקורלציות בין \asset{SPY} לנכסי סיכון אחרים משתנות לפי המצב. שתי התוצאות הן תיאוריות בלבד, אך הן מראות שהמצב החבוי קשור למאפיינים רחבים יותר של השוק ולא רק לאחד המשתנים שהוזנו ישירות למודל.

מכאן מגיע הביטוי \tech{conditional market context}. בעברית, הכוונה היא שה־HMM מספק הקשר מותנה למצב השוק: במקום לשאול רק ``האם מחר תהיה עלייה?'', אפשר לשאול ``באיזה סוג סביבה אנחנו נמצאים עכשיו, כמה המודל בטוח בכך, ועד כמה ההתנהגות של נכסים אחרים משתנה בסביבה הזאת?''. זו יכולת שאינה נמדדת ישירות באמצעות DPA.

\subsection{עושר המודל מול יציבות}

מספר גדול יותר של מצבים מאפשר למודל לתאר מבנה עשיר יותר, אך הוא גם מגדיל את הסיכון למצבים נדירים או לפתרונות שתלויים באתחול. זהו מתח טבעי בין גמישות לבין יציבות.

בממצאים שלנו, \(K=2\) נוטה להיות יציב מאוד אך גס יחסית. \(K=4\) מאפשר הפרדה מפורטת יותר, וב־\asset{SPY}, \asset{QQQ} ו־\asset{NVDA} הוא גם יציב מאוד בין אתחולים. לעומת זאת, ב־\asset{BTC-USD} וב־\asset{GLD} היציבות נמוכה יותר. לכן אין הצדקה לקבוע מראש שמספר מצבים אחד מתאים לכל נכס או לכל תקופה.

\subsection{מדוע לא קראנו למצבים שוק שורי ושוק דובי (\tech{Bull/Bear})?}

בשפה פיננסית, \tech{bull market} הוא בדרך כלל שוק בעל מגמת עלייה מתמשכת, ואילו \tech{bear market} מתאר תקופה ממושכת של ירידות. אלה מושגים כלכליים מוכרים, אך ה־HMM שלנו לא אומן לזהות אותם באופן ישיר.

המצבים בפרויקט נלמדים מתוך ארבע תכונות יומיות, וההבדלים הבולטים ביניהם קשורים בעיקר לתנודתיות, לטווח היומי, ל־\tech{drawdown} ולהתמדה. מצב יכול להיות תנודתי מאוד ועדיין לקבל תשואה ממוצעת חיובית, ולכן תווית ``Bear'' עלולה להיות מטעה. באותו אופן, מספר מצב כמו 0 או 3 הוא שרירותי לחלוטין ויכול להתחלף בין ריצות.

לכן נעשה שימוש בשמות תיאוריים זהירים כמו \tech{Calm / low-vol} או \tech{Rare stress / crisis-like} רק לאחר בחינת הסטטיסטיקות. לדוגמה, מצב 2 של \asset{SPY} מכונה דמוי־משבר משום שהוא נדיר, בעל תשואה ממוצעת שלילית, תנודתיות גבוהה ו־\tech{drawdown} עמוק. גם שם זה הוא פרשנות ולא אמת חיצונית שניתנה למודל.

\subsection{משך המשטר ומה משמעות ה־HMM במקרה שלנו?}

ההשוואה בין משך השהייה התיאורטי, \(1/(1-a_{ii})\), לבין משך השהייה שנמדד בפועל לא חשפה פער גדול בממוצע עבור \asset{SPY}. לכן אין לנו ראיה לכך שהנחת המשך הגאומטרי נכשלת בניסוי הנוכחי. עם זאת, בדיקת ממוצע בלבד אינה בוחנת זנבות ארוכים או כמה סוגים שונים של משכי שהייה; מודל \tech{Hidden Semi-Markov} יכול להיות הרחבה עתידית אם רוצים למדל משך באופן מפורש יותר \cite{bulla2006application}.

המשמעות הכוללת של התוצאות היא שה־\tech{Gaussian HMM} מתאים כאן יותר ל־\textbf{אמידת מצב חבוי} מאשר ל־\textbf{חיזוי נקודתי}. הביטוי \tech{latent state estimation} פירושו שהמודל מנסה להעריך באיזה מצב בלתי־נצפה נמצאת המערכת כרגע. הביטוי \tech{point forecasting} פירושו ניסיון לתת תחזית מספרית או כיוונית לערך הבא בסדרה.

במילים פשוטות, ה־HMM לא הצטיין בשאלה ``כמה תהיה התשואה מחר או מה יהיה הסימן שלה?''. הוא היה שימושי יותר בשאלה ``איזה סוג שוק אנחנו רואים כרגע, כמה זמן מצב כזה נוטה להימשך, ומה מאפיין אותו?''. זו אינה נסיגה מהמטרה המקורית, אלא מסקנה אמפירית שמחדדת היכן המודל מספק ערך בפרויקט זה.
